% ============================================================
% frontmatter/abstract.tex
% ============================================================

\chapter*{Abstract}
\addcontentsline{toc}{chapter}{Abstract}
\noindent

The distribution of neutral hydrogen and its connection to the underlying
dark matter field is a powerful cosmological proben to study the large scale structure and cosmic evolution. Predicting observable
signals, however, is hindered by the large computational cost of
hydrodynamical simulations. This thesis develops a suite of machine
learning methods, based on generative diffusion models to replace
or augment these simulations at field level.

We first introduce \LODI, a pipeline combining an attention-based ResUNet
(\HALOgen) with a conditional variational diffusion model to generate
high-fidelity 21\,cm brightness temperature maps from dark matter-only
simulations, recovering the power spectrum to within 10\% for
$k \leq 10\hMpci$.

We then present \FOLD, which extends this approach to generate
arbitrarily large 3D cosmological fields via a novel sliding window technique, enabling generalization from $25^3\Mpch$ training volumes to
\TNG300-scale fields while preserving large-scale structure and improving high order correlations beyond the 2 point correlations.
We also apply the same \FOLD technique to generate field level posteriors, while being constrained by Planck priors.   

Next, we implement Diffusion Posterior Sampling (DPS) for
field level line of sight reconstruction of the full physical state of the intergalactic medium from Lyman-$\alpha$ forest sightlines using a differentiable Voigt profile forward model.

Finally, we develop a fast neural emulator mapping dark matter fields to HI
and galaxy fields, conditioned on six astrophysical and cosmological
parameters, and embed it within a simulation-based inference pipeline to
constrain $\Om$ and $\sige$ from field-level summaries, while quantifying the impact of field level multi field inference.